\documentclass[11pt, a4paper]{article}

% ── Packages ────────────────────────────────────────────────────────────────
\usepackage[a4paper, margin=2.5cm]{geometry}
\usepackage{amsmath, amssymb, amsthm}
\usepackage{mathtools}
\usepackage{booktabs}
\usepackage{graphicx}
\usepackage{caption}
\usepackage{subcaption}
\usepackage{float}
\usepackage{hyperref}
\usepackage{xcolor}
\usepackage{microtype}
\usepackage{parskip}
\usepackage{enumitem}
\usepackage[T1]{fontenc}
\usepackage{lmodern}

% ── Hyperlink style ──────────────────────────────────────────────────────────
\hypersetup{
  colorlinks = true,
  linkcolor  = blue!60!black,
  citecolor  = blue!60!black,
  urlcolor   = blue!60!black,
}

% ── Theorem environments ─────────────────────────────────────────────────────
\newtheorem{remark}{Remark}

% ── Custom commands ──────────────────────────────────────────────────────────
\newcommand{\R}{\mathbb{R}}
\newcommand{\E}{\mathbb{E}}
\newcommand{\norm}[1]{\left\lVert#1\right\rVert}
\newcommand{\T}{^\top}

% ── Figure path ──────────────────────────────────────────────────────────────
\graphicspath{{results/figures/}}

% ─────────────────────────────────────────────────────────────────────────────
\begin{document}

% ── Title ────────────────────────────────────────────────────────────────────
\begin{center}
  {\LARGE \textbf{LQR vs Reinforcement Learning on a}}\\[4pt]
  {\LARGE \textbf{Discrete-Time Double Integrator}}\\[12pt]
  {\large Ojas Phadake}\\[4pt]
  {\normalsize \textit{Personal Portfolio Project} \quad|\quad \today}
\end{center}

\vspace{4pt}
\hrule
\vspace{14pt}

% ── Abstract ─────────────────────────────────────────────────────────────────
\begin{abstract}
We compare Linear Quadratic Regulation (LQR) and Proximal Policy
Optimisation (PPO) on a discrete-time double integrator, a canonical
second-order system in classical control theory. LQR solves the
infinite-horizon quadratic cost problem analytically by computing the
solution to the Discrete Algebraic Riccati Equation (DARE), yielding an
optimal linear feedback law $u_k = -Kx_k$. PPO, a model-free
reinforcement learning algorithm, is given only the reward signal $r_k =
-(x_k\T Q x_k + u_k\T R u_k)$ and no knowledge of the system dynamics.
After $10^5$ environment interactions, PPO achieves a total cost of
$365.05$, compared to the LQR cost of $364.09$ — a suboptimality gap of
$+0.26\%$. The PPO control signal correlates with the LQR signal at $r =
0.987$, indicating that the learned policy closely approximates the linear
feedback law $u \approx -Kx$ without ever observing the system matrices
$A$ or $B$. This result illustrates the deepest connection between optimal
control and reinforcement learning: both are solving the same Bellman
equation, using different information.
\end{abstract}

\vspace{8pt}
\hrule
\vspace{12pt}

% ─────────────────────────────────────────────────────────────────────────────
\section{Problem Formulation}
% ─────────────────────────────────────────────────────────────────────────────

\subsection{System Dynamics}

We consider the discrete-time double integrator with state $x_k =
[p_k,\; v_k]\T \in \R^2$, where $p_k$ denotes position and $v_k$
velocity, driven by a scalar control input $u_k \in \R$:

\begin{equation}
  x_{k+1} = Ax_k + Bu_k, \qquad
  A = \begin{bmatrix}1 & 0.1 \\ 0 & 1\end{bmatrix}, \quad
  B = \begin{bmatrix}0 \\ 0.1\end{bmatrix}
  \label{eq:dynamics}
\end{equation}

This corresponds to Euler discretisation of the continuous double
integrator $\ddot{p} = u$ with timestep $\Delta t = 0.1\,\mathrm{s}$.
The initial condition is $x_0 = [5,\;0]\T$ and the control objective
is to drive $x_k \to 0$.

The open-loop system has both eigenvalues of $A$ at $+1$ (on the unit
circle), making it marginally stable but not asymptotically stable. No
passive convergence to the origin occurs; an active feedback controller
is required. The system is fully state-controllable, which can be verified
by checking that $\operatorname{rank}([B,\; AB]) = 2$.

\subsection{Shared Objective}

Both LQR and PPO minimise (respectively maximise) the same quadratic
objective. The infinite-horizon cost to be minimised is

\begin{equation}
  J = \sum_{k=0}^{\infty}
      \left( x_k\T Q x_k + u_k\T R u_k \right),
  \qquad
  Q = I_2, \quad R = 0.1
  \label{eq:cost}
\end{equation}

$Q = I_2$ penalises position and velocity deviations equally; $R = 0.1$
permits relatively large control inputs to achieve fast settling.

The RL reward at each timestep is defined as the negative stage cost:

\begin{equation}
  r_k = -\!\left( x_k\T Q x_k + u_k\T R u_k \right)
  \label{eq:reward}
\end{equation}

Maximising $\sum_{k} r_k$ is therefore \emph{identical} to minimising
$J$. This equivalence is not a modelling choice — it is the
mathematical bridge between optimal control and reinforcement learning,
and it is the central observation of this report.

% ─────────────────────────────────────────────────────────────────────────────
\section{Theoretical Background}
% ─────────────────────────────────────────────────────────────────────────────

\subsection{Bellman's Principle of Optimality}

Define the cost-to-go (value function) from state $x$ at time $k$ as

\begin{equation}
  V^*(x_k) = \min_{u_k, u_{k+1},\ldots}
              \sum_{j=k}^{\infty}
              \left( x_j\T Q x_j + u_j\T R u_j \right)
  \label{eq:V}
\end{equation}

Bellman's principle states that an optimal policy satisfies the
fixed-point equation

\begin{equation}
  V^*(x) = \min_{u}
            \left[ x\T Qx + u\T Ru + V^*(Ax + Bu) \right]
  \label{eq:bellman}
\end{equation}

This is the \textit{discrete-time Bellman optimality equation}. It says:
the optimal total cost from $x$ equals the immediate stage cost plus the
optimal cost-to-go from the next state. Every algorithm in this report
is, at its core, attempting to solve or approximate equation
\eqref{eq:bellman}.

\subsection{Dynamic Programming and the Riccati Equation}

For the linear-quadratic problem, the value function has the quadratic
form $V^*(x) = x\T P x$ for some positive-definite matrix $P \in
\R^{2\times 2}$. Substituting into \eqref{eq:bellman} and optimising
over $u$:

\begin{equation}
  x\T Px = \min_{u}
           \left[ x\T Qx + u\T Ru
                  + (Ax+Bu)\T P(Ax+Bu) \right]
\end{equation}

Taking $\partial/\partial u = 0$ gives the optimal control law:

\begin{equation}
  u^* = -\underbrace{(R + B\T PB)^{-1} B\T PA}_{K}\, x
  \label{eq:lqr_law}
\end{equation}

Substituting $u^*$ back yields the \textbf{Discrete Algebraic Riccati
Equation (DARE)}:

\begin{equation}
  P = Q + A\T PA - A\T PB(R + B\T PB)^{-1}B\T PA
  \label{eq:dare}
\end{equation}

This is the fixed-point equation for $P$. Its unique positive-definite
solution yields the globally optimal infinite-horizon controller. The
matrix $P$ \emph{is} the value function: $V^*(x) = x\T Px$.

\begin{remark}
  The DARE \eqref{eq:dare} and the RL Bellman equation
  \eqref{eq:bellman} are the same equation. In LQR we know $A$ and $B$,
  so we can solve for $P$ analytically. In RL we do not know $A$ or $B$,
  so we approximate $V^*$ (or equivalently $P$) from interaction data.
\end{remark}

\subsection{RL as Approximate Dynamic Programming}

In the RL framework, the agent does not observe $A$ or $B$. It observes
transitions $(x_k, u_k, r_k, x_{k+1})$ collected by executing its
policy in the environment. The goal is to find a policy $\pi^*(x)$ such
that $\sum_k r_k$ is maximised — which, by construction of \eqref{eq:reward},
coincides with minimising $J$.

The PPO algorithm maintains two function approximators:

\begin{itemize}[leftmargin=2em]
  \item \textbf{Actor} $\pi_\theta(u \mid x)$: a Gaussian policy
        parameterised by a neural network with weights $\theta$. For
        continuous actions, the network outputs the mean $\mu_\theta(x)$
        and a learned log-standard-deviation.

  \item \textbf{Critic} $V_\phi(x)$: a scalar value function
        approximator parameterised by weights $\phi$. It approximates
        $V^*(x)$, i.e.\ the neural network is learning to represent
        $x\T Px$ for the unknown $P$.
\end{itemize}

PPO updates both networks by collecting rollouts of $T$ steps, computing
Generalised Advantage Estimates (GAE):

\begin{equation}
  \hat{A}_t = \sum_{l=0}^{T-t-1} (\gamma\lambda)^l \delta_{t+l},
  \qquad
  \delta_t = r_t + \gamma V_\phi(x_{t+1}) - V_\phi(x_t)
  \label{eq:gae}
\end{equation}

and then maximising the clipped surrogate objective for $K$ epochs:

\begin{equation}
  \mathcal{L}(\theta) =
  \E_t\!\left[
    \min\!\left(
      \rho_t(\theta)\,\hat{A}_t,\;
      \operatorname{clip}(\rho_t(\theta),\,1-\varepsilon,\,1+\varepsilon)\,\hat{A}_t
    \right)
  \right]
  \label{eq:ppo}
\end{equation}

where $\rho_t(\theta) = \pi_\theta(u_t\mid x_t)/\pi_{\theta_{\text{old}}}(u_t\mid x_t)$
is the probability ratio between the updated and the previous policy, and
$\varepsilon = 0.2$ is the clip threshold. The clip prevents destructively
large policy updates by refusing to credit changes in $\rho_t$ beyond
$[1-\varepsilon,\, 1+\varepsilon]$.

% ─────────────────────────────────────────────────────────────────────────────
\section{Implementation}
% ─────────────────────────────────────────────────────────────────────────────

\subsection{Shared Environment}

Both controllers are evaluated using the same \texttt{DoubleIntegratorEnv}
Gymnasium environment to ensure a fair comparison. The
\texttt{env.simulate(controller\_fn)} method runs one episode from
$x_0 = [5, 0]\T$, collecting state, action, reward, and cost trajectories.
The LQR controller calls \texttt{simulate} directly with $u_k = -Kx_k$;
PPO's trained policy is wrapped in the same signature.

\subsection{LQR Implementation}

The DARE \eqref{eq:dare} is solved using
\texttt{scipy.linalg.solve\_discrete\_are}, which implements a
Schur-decomposition-based algorithm. The resulting matrices are:

\begin{equation}
  P = \begin{bmatrix} 13.827 & 3.868 \\ 3.868 & 4.962 \end{bmatrix},
  \qquad
  K = \begin{bmatrix} 2.585 & 3.575 \end{bmatrix}
  \label{eq:solution}
\end{equation}

giving the feedback law $u_k = -2.585\,p_k - 3.575\,v_k$. The
closed-loop eigenvalues of $A - BK$ are $\{0.899,\; 0.743\}$, both
strictly inside the unit circle, confirming asymptotic stability.

The theoretical minimum cost from $x_0 = [5,\;0]\T$ is

\begin{equation}
  J^* = x_0\T P x_0
  = \begin{bmatrix}5 & 0\end{bmatrix}
    \begin{bmatrix}13.827 & 3.868 \\ 3.868 & 4.962\end{bmatrix}
    \begin{bmatrix}5 \\ 0\end{bmatrix}
  = 345.68
  \label{eq:Jstar}
\end{equation}

This is a hard lower bound. No causal controller, linear or nonlinear,
can achieve $J < J^*$ for this cost function and initial condition.

\subsection{PPO Implementation}

The PPO agent uses Stable-Baselines3 with the following configuration:
$4$ parallel training environments, $n_{\text{steps}} = 1024$ steps per
environment per rollout (giving $4096$ transitions per update),
$10$ optimisation epochs per rollout, minibatch size $64$, learning rate
$3\times10^{-4}$, clip range $\varepsilon=0.2$, and GAE parameter
$\lambda=0.95$. Both the actor and critic are two-layer MLPs with $64$
hidden units and \texttt{tanh} activations.

A \texttt{RescaleAction} wrapper maps the environment's action space
$[-5,5]$ to $[-1,1]$ for the agent, which is standard practice for
Gaussian policies in Stable-Baselines3. The reward is used without
normalisation, preserving its direct correspondence with the LQR cost.

Training runs for $10^5$ total timesteps on a single CPU in
approximately 2--3 minutes. The policy is evaluated deterministically
(mean of the Gaussian, no sampling) on the raw unwrapped environment
to produce a trajectory directly comparable to LQR.

% ─────────────────────────────────────────────────────────────────────────────
\section{Results}
% ─────────────────────────────────────────────────────────────────────────────

\subsection{Open-Loop Dynamics}

Before any controller is applied, it is instructive to observe the
uncontrolled system. With $u_k = 0$, the dynamics reduce to
$x_{k+1} = Ax_k$, and since $A = [[1,\,0.1];\,[0,\,1]]$, the position
remains constant at $p_0 = 5\,\mathrm{m}$ while velocity stays at
zero. This reflects the marginal stability of the open-loop system.
A naive sign controller $u_k = -\operatorname{sign}(p_k)$, which ignores
velocity, produces sustained oscillations due to the lack of damping.
Only full-state feedback — using both $p_k$ and $v_k$ — achieves
convergence.

\subsection{LQR Trajectory}

\begin{figure}[H]
  \centering
  \includegraphics[width=0.72\textwidth]{lqr_trajectory.png}
  \caption{LQR closed-loop trajectory from $x_0 = [5,\;0]\T$.
           \textit{Top}: position converges to zero with slight overshoot.
           \textit{Middle}: velocity rises briefly as the controller
           accelerates the system, then decays.
           \textit{Bottom}: initial control saturates at $u=-5$
           (actuator limit), then smoothly decays to zero.}
  \label{fig:lqr}
\end{figure}

The LQR controller settles to $|p_k| < 0.1\,\mathrm{m}$ in $4.40\,\mathrm{s}$
(44 timesteps) and reaches $p_k \approx 0$ with zero residual error.
The initial control input saturates at the actuator limit $u = -5$ because
the unconstrained optimal input $u_0^* = -Kx_0 = -2.585 \times 5 = -12.93$
exceeds the bound. This saturation introduces a small suboptimality gap
relative to $J^*$, as the DARE solution was derived without explicit input
constraints.

\subsection{PPO Trajectory}

\begin{figure}[H]
  \centering
  \includegraphics[width=0.72\textwidth]{ppo_trajectory.png}
  \caption{PPO closed-loop trajectory from $x_0 = [5,\;0]\T$ after
           $10^5$ training timesteps. The trajectory is qualitatively
           identical to Figure~\ref{fig:lqr}: same settling time,
           same saturation behaviour, negligible residual error.}
  \label{fig:ppo}
\end{figure}

\subsection{Learning Curve}

\begin{figure}[H]
  \centering
  \includegraphics[width=\textwidth]{ppo_learning_curve.png}
  \caption{PPO mean episode reward vs training timesteps. The dashed
           blue line marks the LQR episode reward of $-364.09$.
           Reward improves rapidly from $-776$ at $10\mathrm{k}$ steps
           to within $0.5\%$ of LQR by $50\mathrm{k}$ steps, then
           plateaus.}
  \label{fig:lc}
\end{figure}

The learning curve (Figure~\ref{fig:lc}) shows the characteristic shape
of PPO on a smooth continuous control problem: rapid initial improvement
as coarse structure in the policy is found, followed by slow convergence
as the policy fine-tunes near the optimum. The LQR reward is reached
to within $0.5\%$ at $50\mathrm{k}$ steps and to within $0.26\%$ at
$100\mathrm{k}$ steps.

\subsection{Side-by-Side Comparison}

\begin{figure}[H]
  \centering
  \includegraphics[width=0.80\textwidth]{comparison_overlay.png}
  \caption{Overlay of LQR (blue, solid) and PPO (orange, dashed)
           trajectories. \textit{Top}: position — the two curves are
           nearly indistinguishable. \textit{Middle}: velocity —
           same qualitative shape, minor deviation in the transient.
           \textit{Bottom}: control input — both saturate at $-5$ on
           the first step; the PPO signal is a close but imperfect
           replica of the LQR signal ($\text{MAE} = 0.040$).}
  \label{fig:overlay}
\end{figure}

\begin{figure}[H]
  \centering
  \includegraphics[width=\textwidth]{cumulative_cost.png}
  \caption{Cumulative stage cost $\sum_{i=0}^{k} c_i$ vs time.
           The LQR curve (blue) is the lower envelope — no causal
           controller with knowledge constraints can achieve less.
           The PPO curve (orange) tracks it closely, converging to
           a final cost $0.26\%$ above LQR.}
  \label{fig:cumcost}
\end{figure}

\subsection{Quantitative Comparison}

\begin{table}[H]
  \centering
  \caption{LQR vs PPO on the discrete-time double integrator
           ($x_0 = [5,\;0]\T$, $Q = I_2$, $R = 0.1$).}
  \label{tab:comparison}
  \begin{tabular}{lcc}
    \toprule
    \textbf{Metric} & \textbf{LQR} & \textbf{PPO} \\
    \midrule
    Total cost $J$
      & $364.09$ & $365.05$ \\
    Theoretical minimum $J^* = x_0\T P x_0$
      & \multicolumn{2}{c}{$345.68$} \\
    Suboptimality vs $J^*$
      & $5.33\%$ & $5.60\%$ \\
    \midrule
    Settling time ($|p_k| < 0.1\,\mathrm{m}$)
      & $4.40\,\mathrm{s}$ & $4.50\,\mathrm{s}$ \\
    Peak $|u|$
      & $5.000$ & $5.000$ \\
    RMS $|u|$
      & $0.911$ & $0.891$ \\
    Final $|p|$ (m)
      & $0.000000$ & $0.006575$ \\
    \midrule
    PPO cost gap vs LQR
      & \multicolumn{2}{c}{$+0.26\%$} \\
    Control signal correlation $r$
      & \multicolumn{2}{c}{$0.987$} \\
    Control signal MAE
      & \multicolumn{2}{c}{$0.040$} \\
    \bottomrule
  \end{tabular}
\end{table}

% ─────────────────────────────────────────────────────────────────────────────
\section{Discussion}
% ─────────────────────────────────────────────────────────────────────────────

\subsection{What the Numbers Show}

The $+0.26\%$ cost gap and $r=0.987$ action correlation together tell a
precise story: PPO has effectively rediscovered the linear feedback law
$u \approx -Kx$ without ever observing the system matrices $A$ or $B$.
The neural network — two layers of 64 units with \texttt{tanh}
activations — has learned to represent what is, for this linear-quadratic
problem, essentially a $1\times2$ linear map.

The settling time difference of $0.10\,\mathrm{s}$ (one timestep) is
not statistically meaningful. Both controllers drive the system to rest
in approximately $4.4\,\mathrm{s}$ from $5\,\mathrm{m}$.

\subsection{The 5.33\% Gap Between LQR and $J^*$}

It is important to understand why the LQR controller itself does not
achieve $J^*$. The theoretical minimum \eqref{eq:Jstar} is derived for
the \emph{unconstrained, infinite-horizon} problem. Two sources of
suboptimality exist in simulation:

\begin{enumerate}[leftmargin=2em]
  \item \textbf{Finite horizon}: The simulation runs for $200$ steps
        ($20\,\mathrm{s}$). The tail of the infinite sum
        $\sum_{k=200}^{\infty} c_k$ is nonzero but negligible (the
        system has settled to machine precision by step $\sim$80).
        This contributes only a fraction of the $5.33\%$ gap.

  \item \textbf{Actuator saturation}: The unconstrained LQR law demands
        $u_0 = -Kx_0 = -12.93$ at the initial step, but the actuator
        is clipped to $[-5, 5]$. During the saturated phase the controller
        is suboptimal with respect to $J^*$. This is the dominant source
        of the gap. Since PPO also saturates ($u_0 = -5$), both
        controllers share this constraint equally.
\end{enumerate}

\subsection{Where LQR and RL Diverge}

The near-identical performance here is a consequence of the problem
structure: the system is linear, the cost is quadratic, and the optimal
policy is a linear function of the state. In this regime, a two-layer
neural network has more than sufficient capacity to represent the optimal
policy, and the environment provides a smooth, well-conditioned reward
signal.

The two approaches diverge significantly as the problem becomes harder:

\begin{itemize}[leftmargin=2em]
  \item \textbf{Nonlinear dynamics}: LQR is not applicable. Classical
        extensions (iLQR, DDP) require either known dynamics or a
        differentiable model. PPO requires no model at all.

  \item \textbf{Unknown dynamics}: LQR requires $A$ and $B$ to solve the
        DARE. PPO requires only a simulator or a real physical system.
        This is the key practical advantage of model-free RL.

  \item \textbf{Sample efficiency}: LQR produces the globally optimal
        controller in microseconds (one DARE solve). PPO required
        $10^5$ environment interactions. For expensive-to-run simulators
        or real hardware, this cost is significant.

  \item \textbf{Constraint handling}: Input constraints ($u \in [-5,5]$)
        are handled approximately by PPO (through clipping it experiences
        during training) but are not accounted for in the DARE formulation.
        Model Predictive Control (MPC) would be the appropriate classical
        extension for constrained problems.
\end{itemize}

\subsection{Mathematical Unity}

The central insight of this project is that LQR and RL are not competing
paradigms — they are two implementations of the same idea. Both seek a
policy $\pi^*$ that satisfies the Bellman optimality equation
\eqref{eq:bellman}. LQR exploits the linear-quadratic structure to obtain
a closed-form solution for the value function $V^*(x) = x\T Px$. PPO
approximates $V^*$ with a neural network, using only sampled transitions.

The matrix $P$ computed by the DARE \emph{is} the value function. The
gain $K = (R + B\T PB)^{-1}B\T PA$ \emph{is} the optimal policy. The
neural network in PPO is approximating both, without knowing that the
true solution is quadratic or that the optimal policy is linear.

% ─────────────────────────────────────────────────────────────────────────────
\section{Conclusion}
% ─────────────────────────────────────────────────────────────────────────────

LQR obtains the analytical optimum for this problem because the system
dynamics are known and the cost is quadratic, allowing the Riccati
equation to be solved exactly. PPO learns a near-optimal control policy
through environment interaction alone, without solving the Riccati
equation or observing the system matrices, achieving a cost within
$0.26\%$ of the LQR baseline after $10^5$ training steps.

The action correlation of $r = 0.987$ confirms that PPO has effectively
discovered the linear feedback structure $u \approx -Kx$ that LQR
computes analytically. For the linear-quadratic setting, the two
controllers are empirically indistinguishable.

This equivalence has a precise mathematical basis: both algorithms are
solving the discrete-time Bellman equation. The difference lies not in
the objective but in the available information. LQR uses knowledge of
$A$ and $B$ to compute an exact solution; PPO uses interaction data to
approximate the same solution. As problems grow nonlinear, high-dimensional,
or involve unknown dynamics, the RL approach becomes the only tractable
option — but the Bellman equation it is solving remains the same.

% ─────────────────────────────────────────────────────────────────────────────
\appendix
\section{Hyperparameter Summary}
% ─────────────────────────────────────────────────────────────────────────────

\begin{table}[H]
  \centering
  \caption{PPO hyperparameters used for all experiments.}
  \label{tab:hp}
  \begin{tabular}{llp{6.5cm}}
    \toprule
    \textbf{Parameter} & \textbf{Value} & \textbf{Role} \\
    \midrule
    \texttt{total\_timesteps} & $10^5$ & Total environment steps across all envs \\
    \texttt{n\_envs}          & $4$    & Parallel training environments \\
    \texttt{n\_steps}         & $1024$ & Steps per env per rollout \\
    \texttt{batch\_size}      & $64$   & Mini-batch size for gradient updates \\
    \texttt{n\_epochs}        & $10$   & Gradient epochs per rollout \\
    \texttt{learning\_rate}   & $3\times10^{-4}$ & Adam optimiser step size \\
    \texttt{clip\_range} $\varepsilon$ & $0.2$ & PPO probability ratio clip \\
    \texttt{gae\_lambda} $\lambda$ & $0.95$ & GAE bias-variance tradeoff \\
    \texttt{gamma} $\gamma$   & $0.99$ & Discount factor \\
    \texttt{ent\_coef}        & $0.0$  & Entropy bonus coefficient \\
    \texttt{vf\_coef}         & $0.5$  & Value loss weight in combined loss \\
    \texttt{net\_arch}        & $[64, 64]$ & Hidden layer sizes (both actor and critic) \\
    \texttt{activation}       & \texttt{tanh} & Nonlinearity \\
    \bottomrule
  \end{tabular}
\end{table}

\section{Code Structure}
\label{app:code}

The full project is structured as follows:

\begin{verbatim}
LQRvsRL/
|-- envs/
|   `-- double_integrator_env.py   # Shared Gymnasium environment
|-- controllers/
|   |-- lqr.py                     # DARE solver, LQRController class
|   `-- rl_agent.py                # PPO training, evaluation, callback
|-- analysis/
|   |-- plots.py                   # All figure generation
|   `-- compare.py                 # Metrics, comparison table, LaTeX output
|-- results/
|   |-- figures/                   # All saved plots
|   `-- data/                      # Saved trajectories and model
|-- run_lqr.py                     # Entry point: dynamics + LQR
|-- train_ppo.py                   # Entry point: PPO training
`-- main_compare.py                # Entry point: full comparison
\end{verbatim}

\end{document}